\documentclass{article}
\usepackage[utf8]{inputenc}
\usepackage[T1]{fontenc}
\usepackage{lmodern}
\usepackage{hyperref}
\usepackage{geometry}
\geometry{a4paper, margin=1in}

\title{Prompt Engineering Quiz: Week 3 \& 4}
\author{Ankara Medipol University}
\date{\today}

\begin{document}

\maketitle

\section*{Instructions}
Please choose the best answer for each of the following questions.

\begin{enumerate}

    \item Which component of a good prompt specifies the persona the AI should adopt and the target reader?
    \begin{enumerate}
        \item Output Spec
        \item Context
        \item Role/Audience
        \item Constraints
    \end{enumerate}

    \item When you provide 1-3 high-quality examples of the desired output format in your prompt, what pattern are you using?
    \begin{enumerate}
        \item Zero-shot prompting
        \item Few-shot prompting
        \item Role prompting
        \item Reflection prompting
    \end{enumerate}

    \item What is the primary purpose of Retrieval-Augmented Generation (RAG)?
    \begin{enumerate}
        \item To make the model's responses more creative and imaginative.
        \item To ground the model's answers in provided, specific context.
        \item To translate the model's output into different languages.
        \item To have the model critique its own work before finalizing an answer.
    \end{enumerate}

    \item The instruction "Let's think step by step" is a common way to trigger which reasoning technique, especially when no examples are provided?
    \begin{enumerate}
        \item Tree-of-Thoughts (ToT)
        \item Zero-shot Chain-of-Thought (CoT)
        \item Self-consistency
        \item Plan-then-solve
    \end{enumerate}

    \item Why are delimiters like triple backticks (\texttt{\`\`\`}) important in a prompt?
    \begin{enumerate}
        \item They make the prompt run faster.
        \item They are required for all prompts to work.
        \item They clearly separate instructions from data or context, reducing ambiguity.
        \item They automatically format the output as a code block.
    \end{enumerate}

    \item Asking the model to "be brief" and "include all details" in the same prompt is an example of which anti-pattern?
    \begin{enumerate}
        \item Overloaded ask
        \item Vague audience
        \item Hidden constraints
        \item Conflicting directives
    \end{enumerate}

    \item Which framework is adapted from project management and uses the components: Context, Objective, Style, Tone, Audience, and Response?
    \begin{enumerate}
        \item SMART
        \item RACI
        \item IDEA
        \item CO-STAR
    \end{enumerate}

    \item What is the main difference between a prompt pattern and a prompt framework?
    \begin{enumerate}
        \item Patterns are for simple tasks, while frameworks are only for complex AI.
        \item Frameworks package multiple patterns into a reusable, structured template for a workflow.
        \item Patterns define the output format, while frameworks define the input data.
        \item There is no difference; the terms are interchangeable.
    \end{enumerate}

    \item The process of using an LLM to generate, critique, or refine another prompt is known as:
    \begin{enumerate}
        \item Prompt Chaining
        \item Meta-Prompting
        \item Dynamic Few-Shot Selection
        \item Pattern Stacking
    \end{enumerate}

    \item In the context of LLMs, what does "Tool Use" or "Function Calling" allow the model to do?
    \begin{enumerate}
        \item Execute Python code directly within the model's environment.
        \item Browse the internet without any restrictions.
        \item Generate a structured JSON object specifying a function to be called and its arguments.
        \item Automatically debug and fix errors in its own code.
    \end{enumerate}

    \item A comprehensive organizational asset that details a complete, multi-step workflow with prompts, human review steps, and guardrails is best described as a:
    \begin{enumerate}
        \item Prompt Pattern
        \item Prompt Framework
        \item Prompt Playbook
        \item Prompt Template
    \end{enumerate}

    \item The IDEA Loop is a framework for continuous improvement. What does the "A" in IDEA stand for?
    \begin{enumerate}
        \item Analyze
        \item Adjust
        \item Audience
        \item Action
    \end{enumerate}

\end{enumerate}

\end{document}
